Infrastructuur management bestaat uit een aantal verschillende onderdelen:
\begin{description}
\item[Deployment] Het opbrengen en installeren van netwerk connected devices
\item[Automation and Orchestration] Het configureren en aanpassen van netwerk connected devices
\item[Monitoring] Het in de gaten houden van netwerk connected devices
\item[Event Management] Het afhandelen van meldingen van netwerk connected devices
\end{description}

Deze hele set aan zaken wordt vaak aangeduid als ITOM (IT Operations Management). ITOM behandeld alles dat te maken heeft met de zaken die te maken hebben van de meestal in een datacentrum staande systemen om te zorgen dat servers, besturingssystemen en netwerken soepel en betrouwbaar draaien.

Een belangrijk onderdeel hieruit is Infrastructure as Code (IaC).

%Proactive Monitoring: Keeping an eye on the health of servers, applications, and networks to catch anomalies before they disrupt business.
%Automation & Orchestration: Streamlining routine configuration, patching, and provisioning tasks to reduce human error and save operational time.
%IT Asset Management (ITAM): Tracking all hardware and software throughout its lifecycle to control costs and ensure compliance.
%Event Management: Filtering out "noise" from IT systems and routing genuine alerts to the correct teams for immediate troubleshooting
%Infrastructure as Code (IaC): Hiermee beheert en configureert u IT-infrastructuur via code in plaats van handmatige processen.Bekende tools: Terraform en Ansible.
%Infrastructuur Monitoring: Biedt real-time inzicht in de gezondheid, prestaties en beveiliging van servers, netwerken en applicaties.Bekende tools: Datadog, Paessler PRTG en Zabbix
%IT Operations Management (ITOM): Software voor het dagelijks beheer van fysieke resources, datacenters en cloud-infrastructuur.
